\input{estilo}
\begin{document}
\begin{figura}\begin{tikzpicture}[pieza/.style={caja, minimum width=26mm, minimum height=8mm},
    chunk/.style={archivo, minimum width=6mm, minimum height=7.5mm, inner sep=0pt, font=\scriptsize},
    local/.style={archivo=segundo, minimum width=6mm, minimum height=7.5mm, inner sep=0pt, font=\scriptsize, fill=segundo!12, text=segundo}]
  % dos nodos iguales, cada uno con un worker y un chunk server con sus chunks
  \foreach \n/\x/\a/\b in {1/0/1/3, 2/4.2/2/4} {
    \draw[marco] (\x,0) rectangle (\x+3.2,-3.5);
    \node[below=1pt] at (\x+1.6,-3.5) {nodo \n};
    \node[pieza] (w\n) at (\x+1.6,-0.8) {worker};
    \node[pieza, minimum height=18mm] (cs\n) at (\x+1.6,-2.4) {};
    \node[font=\small] at (\x+1.6,-1.8) {chunk server};
    \node[local] (c\n) at (\x+1.15,-2.75) {\a};
    \node[chunk] at (\x+2.05,-2.75) {\b};
  }
  % el coordinador le da a cada nodo el map cuyo input ya tiene en su disco
  \node[nodo, minimum size=13mm] (co) at (3.7,1.9) {};
  \node[left=3pt] at (co.west) {coordinador};
  \draw[->] (co.south west) to[out=225, in=90] node[pos=0.55, etiq, left=2pt, text=segundo] {map 1} (w1.north);
  \draw[->] (co.south east) to[out=-45, in=90] node[pos=0.55, etiq, right=2pt, text=segundo] {map 2} (w2.north);
\end{tikzpicture}\end{figura}
\end{document}
